\section{Conclusion}
\label{sec:conclusion}
Schubert closed his 2007 note wondering whether h\textsubscript{2} of institutes would prove more interesting than h\textsubscript{2} of journals, absent the researcher--institute--country data needed to check. Nineteen years and 30.0 million OpenAlex authors later, the answer is yes, but not in the way a single top-10 table would suggest.

The overall h\textsubscript{2} ranking is the least interesting object in this paper. Because Medicine alone accounts for 25\% of the author pool, the unrestricted university-wide aggregate is, in substance, a ranking of medical school size. The field decomposition reveals a completely different map of global research strength: Wageningen ranks first in Agricultural \& Biological Sciences and Environmental Science by field h\textsubscript{2} while ranking 172nd overall, though the Environmental Science lead holds only under the paper's single-field assignment and not under multi-field assignment (Subsection~\ref{subsec:results-field-robustness}); Carnegie Mellon ranks first in Computer Science by field h\textsubscript{2} from 175th, \textcolor{red}{though, like the Environmental Science lead, this too holds only under the paper's single-field assignment: Carnegie Mellon's Computer Science lead falls to Stanford under the more aggressive 20\% multi-field threshold (Subsection~\ref{subsec:results-field-robustness})}; the University of Science and Technology of China ranks first in Materials Science, Engineering, and Energy from 39th; and Shenyang Pharmaceutical University ranks first in Pharmacology from 535th. Brazil offers perhaps the sharpest illustration of what the overall ranking hides: the Universidade de S\~{a}o Paulo and the Universidade Estadual de Campinas rank among the top six globally in Dentistry by field h\textsubscript{2} despite neither country nor either institution appearing anywhere near the top of any general-purpose ranking. The field-level h\textsubscript{2} is the unit that should be the default for comparison, not the aggregate.

On the methodological side, Egghe's power-law functional form is empirically well-supported: \textcolor{red}{author count alone explains 91\% of the variance in h\textsubscript{2}, and institution count alone explains 92\% of the variance in h\textsubscript{3}}. However, the theoretical compounding prediction $\alpha_0\alpha_1 \approx \beta_1$ fails by 177\%. % The source of the discrepancy is that the discrete MLE tail fits locate $x_\mathrm{min}$ far above unity and describe the extreme upper tail, not the global Lotka parameter Egghe's derivation requires.
The model is empirically useful; the underlying theory needs a different reading. \textcolor{red}{Two external comparisons bear on whether the efficiency normalisation retains a genuine research-quality signal. Against the CWTS Leiden Ranking Open Edition 2024~\citep{leiden2024}, a ranking that shares OpenAlex's underlying data infrastructure and so functions as a convergent check rather than an independent validation, both h\textsubscript{2} and $\varepsilon_2$ correlate with the Leiden mean normalised citation score across 1,495 matched institutions, but raw h\textsubscript{2} ($\rho=0.634$) correlates more strongly than $\varepsilon_2$ ($\rho=0.580$). Against Scimago Institutions Rankings 2026~\citep{scimago2026}, a Scopus-derived ranking independent of OpenAlex, the same pattern recurs and is more pronounced: raw h\textsubscript{2} correlates at $\rho=0.752$ against SIR's Global Rank versus $\rho=0.338$ for $\varepsilon_2$, on a smaller, name-matched set of 3,659 institutions (Subsections~\ref{subsec:methods-leiden}--\ref{subsec:methods-scimago}, Subsubsections~\ref{subsubsec:leiden-validation}--\ref{subsubsec:scimago-validation}). Both comparisons show raw, scale-dominated h\textsubscript{2} tracking external rankings more closely than the size-normalised $\varepsilon_2$ does, plausibly because those external composites are themselves influenced by scale-correlated factors that dividing out size removes; this is reported as a description of the metrics' behaviour, not as validation of the efficiency normalisation's superiority.}

Several limitations constrain interpretation. Author disambiguation is the most consequential open problem: the four outlier cases identified in cross-validation (notably Lei Jiang, $\mathrm{h}_1^\mathrm{OpenAlex} = 215$ against Web of Science $= 115$) confirm that merger errors cluster in non-Anglophone and common-name strata, exactly the authors whose inflated h\textsubscript{1} values would propagate most damagingly into institutional and country-level rankings. That cross-validation sample was deliberately adversarial, however, and cannot by itself say how much of the full \textcolor{red}{30.0-million}-author population is affected; a Monte Carlo sensitivity analysis run at population-wide error rates from 1\% to a deliberately implausible 80\% (Subsections~\ref{subsec:methods-sensitivity}, \ref{subsec:results-sensitivity}) finds the risk bounded rather than fatal, with institution- and country-level rank correlations declining only modestly before plateauing ($\rho \geq 0.988$) rather than degrading toward noise at any rate tested, and the paper's country-comparison claims holding throughout. The field-leadership claims are noisier and should be read with correspondingly more caution than the paper's other claims. Arts and Humanities are systematically under-covered in all major citation databases: both threshold authors sampled at Oxford returned no value in Web of Science, Scopus, or Google Scholar, and h-index comparisons in book-based fields should be read with correspondingly less confidence. Field-specific rankings are also sensitive to how field membership is defined: a multi-field assignment (authors counted in every field comprising $\ge 33\%$ of their output) yields highly stable overall rankings ($\rho = 0.963$) but shifts Wageningen's Environmental Science lead to the University of Colorado Boulder, confirming that the field-leadership claims should be read as reflecting depth of specialisation rather than unconditional supremacy (Subsection~\ref{subsec:results-field-robustness}). OpenAlex's inconsistent retraction flagging leaves a small, unquantified fraction of h\textsubscript{1} values potentially inflated. \textcolor{red}{The top of the small-institution efficiency table (Table~\ref{tab:h2-efficiency-all}) is similarly affected: several entries, including Bellevue University and three other primarily-teaching institutions, owe their apparent efficiency to a handful of author records that are almost certainly disambiguation-merge artifacts rather than genuine researchers (see the discussion following Table~\ref{tab:h2-efficiency-all}); this correction was not extended to the rest of the dataset and should be read as a targeted caveat, not a general audit.} Saudi Arabia's efficiency rank (4th, $\varepsilon_3 = 6.69$) is plausibly an artifact of affiliation-listing for funding rather than resident research presence~\citep{saudi}, a data-quality distortion visible in the output but not removable without ground-truth affiliation data\textcolor{red}{; a multi-institution robustness check (Subsection~\ref{subsec:results-institution-robustness}) rules out the current-affiliation tiebreak rule specifically as the cause ($\varepsilon_3$ rises rather than falls under that check), but Saudi Arabia is simply the case where this kind of distortion happens to be visible, not necessarily the only country affected by it}. \textcolor{red}{Current-affiliation attribution is itself a related, broader limitation: every author's entire career h\textsubscript{1} is credited to their most recent educational institution, so internationally mobile or multiply affiliated researchers can have career-long impact attributed to wherever they are affiliated now rather than to where the underlying work was produced. A multi-institution robustness check (Subsection~\ref{subsec:results-institution-robustness}) bounds one specific piece of this problem, the arbitrary tiebreak applied when two or more institutions are equally most-recent, and finds it moves institution-level rankings only modestly ($\rho=0.966$) and country-level rankings very little ($\rho=0.997$); it does not establish that current-affiliation attribution more broadly is unproblematic, since it leaves untested every author whose single most-recent affiliation is unambiguous but still not where their cited work was conducted. That broader concern is addressed in this paper by narrowing its claims to what current-affiliation-based h\textsubscript{2} and h\textsubscript{3} can support, not by a further re-attribution that would resolve it.} Finally, h\textsubscript{2} is not comparable across fields of different size: Decision Sciences (max h\textsubscript{2} $= 22$) cannot be placed on the same scale as Medicine (max h\textsubscript{2} $= 119$), and the results should be read within fields only.

The most impactful near-term improvement to this framework would require a data source that does not yet exist: departmental affiliation in OpenAlex. If OpenAlex were to record the department within an institution alongside the institution itself, the hierarchy would extend naturally to researcher $\to$ department $\to$ institution $\to$ country, with successive indices h\textsubscript{1} through h\textsubscript{4}. Departments are field-specialized by construction, so the step of inferring each author's modal field from topic weights may become unnecessary at the departmental level. I flag this as the single most valuable addition OpenAlex could make for scientometric work of this kind.